\section{Prediction-Interval Calibration and Tests}
\label{app:intervals}
\setcounter{table}{0}

\subsection{Exact prior-only calibration algorithm}

The confirmatory H4 analysis uses the frozen stage-policy point forecast for
each target and forecast stage.  Historical calibration errors are pooled
only within the same target series, forecast stage, and frozen selected
model.  For a current forecast dated \(t\), a historical error is admissible
only when its stored outcome release date is strictly earlier than \(t\).
The admissible history is sorted by outcome release date, target-period
order, and forecast date; the most recent observations up to the
domain-specific window length are then retained.

Let \(v_j=|e_j|\) denote absolute calibration errors, with \(j=0\) the newest
observation after the admissibility and window rules have been applied.
The 80\% intervals are symmetric around the point forecast.

For GDP the primary half-width is the unweighted empirical 0.80 quantile of
the most recent 20 absolute errors, subject to at least 12 available errors.
The implementation uses NumPy's \texttt{quantile} with
\texttt{method="higher"}; equivalently, the selected half-width is the first
observed order statistic at or above the requested empirical quantile.

For inflation and labour the primary half-width is a weighted empirical
0.80 quantile over at most 48 errors, subject to at least 24 available
errors.  Inflation uses
\[
    w_j = 0.5^{j/18},
\]
while labour uses
\[
    w_j = 0.94^j.
\]
To evaluate the weighted quantile, the absolute errors are stably sorted in
ascending order together with their positive weights.  If
\(v_{(1)}\leq\cdots\leq v_{(M)}\) are the sorted errors, the reported
half-width is
\[
    q_{0.8}^{w}
    =
    v_{(r)},\qquad
    r=\min\left\{
        m:\sum_{i=1}^{m}w_{(i)}
        \geq 0.8\sum_{i=1}^{M}w_{(i)}
    \right\}.
\]
There is no fallback rule.  A forecast with fewer than the required number
of admissible prior errors is H4-ineligible.  The confirmatory evaluation is
stage specific and does not require a balanced set of stages at the end of
the sample.

\subsection{Benchmark interval methods}

GDP is evaluated against one benchmark, a Gaussian RMSE interval.  Inflation
and labour are evaluated against both the Gaussian RMSE interval and an
unweighted rolling empirical \(q_{0.8}\) interval.

The rolling benchmark uses the same 48-observation maximum window and
24-error minimum history as the inflation and labour primary methods, and
uses the same ``higher'' 0.80 empirical quantile convention.  The Gaussian
benchmark uses the same admissible calibration window as the primary method.
Its half-width is
\[
    q^G
    =
    \Phi^{-1}(0.90)
    \sqrt{\frac{1}{M}\sum_{j=1}^{M} e_j^2},
\]
and is centred on the point forecast.  The Gaussian benchmark is therefore
an RMSE-scaled central 80\% interval and does not apply a separate
forecast-bias adjustment.

\subsection{Coverage and independence diagnostics}

Let \(I_t=1\) denote an interval violation and \(I_t=0\) a covered
realisation.  With \(x=\sum_t I_t\), \(N\) observations, and nominal
violation probability \(\alpha=0.20\), the unconditional-coverage statistic
is
\[
    LR_{\mathrm{uc}}
    =
    -2\left[
        \ell(\alpha)-\ell(\widehat p)
    \right],
    \qquad
    \widehat p=x/N,
\]
where
\[
    \ell(p)=x\log p+(N-x)\log(1-p).
\]
Under the reference asymptotics,
\(LR_{\mathrm{uc}}\sim\chi^2_1\).  A 95\% Wilson binomial confidence interval
is also reported for the coverage probability itself.

For violation independence, let \(n_{ij}\) count transitions from state
\(i\) to state \(j\), \(i,j\in\{0,1\}\).  Define
\[
    \widehat\pi
    =
    \frac{n_{01}+n_{11}}{n_{00}+n_{01}+n_{10}+n_{11}},
    \quad
    \widehat\pi_{01}
    =
    \frac{n_{01}}{n_{00}+n_{01}},
    \quad
    \widehat\pi_{11}
    =
    \frac{n_{11}}{n_{10}+n_{11}}.
\]
The independence statistic is
\[
    LR_{\mathrm{ind}}
    =
    -2\left[
        \ell_{\mathrm{iid}}-\ell_{\mathrm{Markov}}
    \right]
    \sim\chi^2_1,
\]
and joint conditional coverage is
\[
    LR_{\mathrm{cc}}
    =
    LR_{\mathrm{uc}}+LR_{\mathrm{ind}}
    \sim\chi^2_2.
\]
In computation, zero-count log-likelihood terms contribute zero; transition
probabilities with a zero denominator are set to zero, matching the
precommitted implementation.  Independence and joint tests are left missing
when a stage cell contains fewer than two interval observations.

\subsection{Sharpness, interval score, and benchmark inference}

Average interval width is
\[
    AIW=N^{-1}\sum_{t=1}^{N}(U_t-L_t).
\]
For \(\alpha=0.20\), the Winkler interval score is
\[
    IS_{\alpha,t}
    =
    (U_t-L_t)
    +\frac{2}{\alpha}(L_t-y_t)\I(y_t<L_t)
    +\frac{2}{\alpha}(y_t-U_t)\I(y_t>U_t).
\]
The primary benchmark comparison is based on the paired score differential
\[
    d_t^{IS}
    =
    IS_t^{\mathrm{primary}}
    -
    IS_t^{\mathrm{benchmark}}.
\]
Negative values favour the primary calibration method.  The reported
standard error of \(\bar d^{IS}\) is Newey--West HAC with Bartlett weights
and bandwidth
\[
    B
    =
    \left\lfloor
        4(N/100)^{2/9}
    \right\rfloor,
\]
capped at \(N-1\).  The two-sided \(p\)-value is based on the corresponding
normal statistic.

\subsection{Frozen calibration hyperparameters}

\begin{table}[H]
\centering
\caption{Frozen H4 interval-calibration specification}
\label{tab:h4_hyperparameters}
\small
\resizebox{0.94\textwidth}{!}{%
\begin{tabular}{lllll}
\toprule
Domain & Primary method & Window & Minimum history & Weight/quantile rule \\
\midrule
GDP & Rolling \(q_{0.8}\) & 20 & 12 & NumPy ``higher'' \\
Inflation & Exp.-weighted \(q_{0.8}\) & 48 & 24 & \(0.5^{j/18}\) \\
Labour & Exp.-weighted \(q_{0.8}\) & 48 & 24 & \(0.94^j\) \\
\bottomrule
\end{tabular}%
}
\par\medskip
\begin{minipage}{0.94\textwidth}
\footnotesize
\textit{Notes:} All methods use only finite historical errors whose outcome
release date is strictly earlier than the current forecast date.  Calibration
is restricted to the same target, forecast stage, and frozen selected model.
There is no fallback below the minimum history.
\end{minipage}
\end{table}
